% --- Archon physics formalization source begin ---
% archon:physics
% archon:covers PhyXMiniProblems/problem_phyx_mini_0951.lean
% archon:source-report reports/phyx_mini/problem_phyx_mini_0951.source.json

\chapter{Physics problem phyx\_mini\_0951}
\label{ch:PhyXMiniProblems_problem_phyx_mini_0951}

\paragraph{Problem source.}
\#\# Physical scenario

A uniform bar has mass \textbackslash{}( 0.0120\textbackslash{},\textbackslash{}text\{kg\} \textbackslash{}) and is \textbackslash{}( 30.0\textbackslash{},\textbackslash{}text\{cm\} \textbackslash{}) long. It pivots without friction about an axis perpendicular to the bar at point \textbackslash{}( a \textbackslash{}). The gravitational force on the bar acts in the \textbackslash{}( -y \textbackslash{})-direction. The bar is in a uniform magnetic field that is directed into the page and has magnitude \textbackslash{}( B = 0.150\textbackslash{},\textbackslash{}text\{T\} \textbackslash{}).

\#\# Question

What must be the current \textbackslash{}( I \textbackslash{}) in the bar for the bar to be in rotational equilibrium when it is at an angle \textbackslash{}( 	heta = 30.0\textasciicircum{}\textbackslash{}circ \textbackslash{}) above the horizontal?

\#\# Answer choices

- A: A: 4.52A

- B: B: 2.26A

- C: C: 15.1A

- D: D: 2.60A

\#\# Figure caption (auxiliary; use the image as primary evidence)

The image shows a diagram with axes labeled \textbackslash{}(x\textbackslash{}) and \textbackslash{}(y\textbackslash{}), forming a coordinate system. There is a diagonal vector, labeled \textbackslash{}(\textbackslash{}vec\{B\}\textbackslash{}), extending from point \textbackslash{}(a\textbackslash{}) on the origin to point \textbackslash{}(b\textbackslash{}). The vector forms an angle \textbackslash{}(\textbackslash{}theta\textbackslash{}) with the x-axis. There are multiple blue crosses in the background, arranged in a grid pattern, likely representing a uniform field perpendicular to the plane of the axes. Points \textbackslash{}(a\textbackslash{}) and \textbackslash{}(b\textbackslash{}) are marked with dots.

\#\# Dataset metadata

Magnetism; Physical Model Grounding Reasoning, Spatial Relation Reasoning

\paragraph{Recorded answer/context.}
B: B: 2.26A

\paragraph{Figure/image path.}
/root/proposal\_for\_physic/science-mango/phyx\_mini\_run/phyx\_data/test\_image/951.png

\paragraph{Formalization target.}
create a compiling Lean file with sorry bodies at `PhyXMiniProblems/problem\_phyx\_mini\_0951.lean`. The Lean declarations must preserve the physical quantities, dimensions or dimensional roles, figure labels, governing-law hypotheses, and final relation expressed by this problem.
Use Mathlib/Physlib names found through LeanExplore where available. If a physics API is missing, introduce faithful local abstractions rather than scalar placeholder aliases.

\begin{theorem}[Physics formalization target]
\label{thm:physics:phyx_mini_0951:target}
\lean{PhyXMiniProblems.ProblemPhyXMini0951.problem_phyx_mini_0951}
\uses{def:physics:phyx-mini-0951:phyxminiproblems-problemphyxmini0951-pivotedconductingbarsetup, def:physics:phyx-mini-0951:phyxminiproblems-problemphyxmini0951-matchesproblemandfigurereadouts, def:physics:phyx-mini-0951:phyxminiproblems-problemphyxmini0951-matchespivoteduniformbarscenario, def:physics:phyx-mini-0951:phyxminiproblems-problemphyxmini0951-hasphysicalbarparameters, def:physics:phyx-mini-0951:phyxminiproblems-problemphyxmini0951-hasuniformappliedmagneticfield, def:physics:phyx-mini-0951:phyxminiproblems-problemphyxmini0951-satisfiesuniformbartorquelaws, def:physics:phyx-mini-0951:phyxminiproblems-problemphyxmini0951-isinrotationalequilibrium, def:physics:phyx-mini-0951:phyxminiproblems-problemphyxmini0951-recordeddatasetanswer, def:physics:phyx-mini-0951:phyxminiproblems-problemphyxmini0951-answermatchescurrenttodisplayedprecision, def:physics:phyx-mini-0951:phyxminiproblems-problemphyxmini0951-isuniqueclosestcurrentchoice}
The assigned autoformalize agent should translate this physics problem into Lean declarations in the covered file, with theorem and lemma proof bodies written as `by sorry`.
\end{theorem}
\begin{proof}
This is an autoformalization task, not a proof task. Produce statements that can later be proved by the physics prover without weakening the physical model.
\end{proof}

% --- Archon declaration topology begin ---
\section{Declaration topology}
\begin{definition}[acceleration Dimension]
  \label{def:physics:phyx-mini-0951:phyxminiproblems-problemphyxmini0951-accelerationdimension}
  \lean{PhyXMiniProblems.ProblemPhyXMini0951.accelerationDimension}
  The physical dimension L T⁻² of acceleration
\end{definition}
\begin{proof}
  This is the declared model definition of acceleration Dimension.
\end{proof}
\begin{definition}[force Dimension]
  \label{def:physics:phyx-mini-0951:phyxminiproblems-problemphyxmini0951-forcedimension}
  \lean{PhyXMiniProblems.ProblemPhyXMini0951.forceDimension}
  The physical dimension M L T⁻² of force
\end{definition}
\begin{proof}
  This is the declared model definition of force Dimension.
\end{proof}
\begin{definition}[torque Dimension]
  \label{def:physics:phyx-mini-0951:phyxminiproblems-problemphyxmini0951-torquedimension}
  \lean{PhyXMiniProblems.ProblemPhyXMini0951.torqueDimension}
  The physical dimension M L² T⁻² of torque
\end{definition}
\begin{proof}
  This is the declared model definition of torque Dimension.
\end{proof}
\begin{definition}[magnetic Flux Density Dimension]
  \label{def:physics:phyx-mini-0951:phyxminiproblems-problemphyxmini0951-magneticfluxdensitydimension}
  \lean{PhyXMiniProblems.ProblemPhyXMini0951.magneticFluxDensityDimension}
  The physical dimension M T⁻¹ C⁻¹ of magnetic flux density
\end{definition}
\begin{proof}
  This is the declared model definition of magnetic Flux Density Dimension.
\end{proof}
\begin{definition}[electric Current Dimension]
  \label{def:physics:phyx-mini-0951:phyxminiproblems-problemphyxmini0951-electriccurrentdimension}
  \lean{PhyXMiniProblems.ProblemPhyXMini0951.electricCurrentDimension}
  The physical dimension C T⁻¹ of electric current
\end{definition}
\begin{proof}
  This is the declared model definition of electric Current Dimension.
\end{proof}
\begin{definition}[Mass Quantity]
  \label{def:physics:phyx-mini-0951:phyxminiproblems-problemphyxmini0951-massquantity}
  \lean{PhyXMiniProblems.ProblemPhyXMini0951.MassQuantity}
  A nonnegative, unit-independent physical mass
\end{definition}
\begin{proof}
  This is the declared model definition of Mass Quantity.
\end{proof}
\begin{definition}[Length Quantity]
  \label{def:physics:phyx-mini-0951:phyxminiproblems-problemphyxmini0951-lengthquantity}
  \lean{PhyXMiniProblems.ProblemPhyXMini0951.LengthQuantity}
  A nonnegative, unit-independent physical length
\end{definition}
\begin{proof}
  This is the declared model definition of Length Quantity.
\end{proof}
\begin{definition}[Acceleration Magnitude]
  \label{def:physics:phyx-mini-0951:phyxminiproblems-problemphyxmini0951-accelerationmagnitude}
  \lean{PhyXMiniProblems.ProblemPhyXMini0951.AccelerationMagnitude}
  \uses{def:physics:phyx-mini-0951:phyxminiproblems-problemphyxmini0951-accelerationdimension}
  A nonnegative, unit-independent acceleration magnitude
\end{definition}
\begin{proof}
  This is the declared model definition of Acceleration Magnitude.
\end{proof}
\begin{definition}[Magnetic Flux Density Magnitude]
  \label{def:physics:phyx-mini-0951:phyxminiproblems-problemphyxmini0951-magneticfluxdensitymagnitude}
  \lean{PhyXMiniProblems.ProblemPhyXMini0951.MagneticFluxDensityMagnitude}
  \uses{def:physics:phyx-mini-0951:phyxminiproblems-problemphyxmini0951-magneticfluxdensitydimension}
  A nonnegative, unit-independent magnetic-flux-density magnitude
\end{definition}
\begin{proof}
  This is the declared model definition of Magnetic Flux Density Magnitude.
\end{proof}
\begin{definition}[Electric Current Quantity]
  \label{def:physics:phyx-mini-0951:phyxminiproblems-problemphyxmini0951-electriccurrentquantity}
  \lean{PhyXMiniProblems.ProblemPhyXMini0951.ElectricCurrentQuantity}
  \uses{def:physics:phyx-mini-0951:phyxminiproblems-problemphyxmini0951-electriccurrentdimension}
  A signed current relative to the orientation from a to b. Its dimensionful representation keeps current distinct from a bare scalar
\end{definition}
\begin{proof}
  This is the declared model definition of Electric Current Quantity.
\end{proof}
\begin{definition}[Force Magnitude]
  \label{def:physics:phyx-mini-0951:phyxminiproblems-problemphyxmini0951-forcemagnitude}
  \lean{PhyXMiniProblems.ProblemPhyXMini0951.ForceMagnitude}
  \uses{def:physics:phyx-mini-0951:phyxminiproblems-problemphyxmini0951-forcedimension}
  A nonnegative, unit-independent force magnitude
\end{definition}
\begin{proof}
  This is the declared model definition of Force Magnitude.
\end{proof}
\begin{definition}[Signed Torque Quantity]
  \label{def:physics:phyx-mini-0951:phyxminiproblems-problemphyxmini0951-signedtorquequantity}
  \lean{PhyXMiniProblems.ProblemPhyXMini0951.SignedTorqueQuantity}
  \uses{def:physics:phyx-mini-0951:phyxminiproblems-problemphyxmini0951-torquedimension}
  A signed torque component about the page-normal pivot axis; positive means counterclockwise in the supplied xy view
\end{definition}
\begin{proof}
  This is the declared model definition of Signed Torque Quantity.
\end{proof}
\begin{definition}[mass Readout]
  \label{def:physics:phyx-mini-0951:phyxminiproblems-problemphyxmini0951-massreadout}
  \lean{PhyXMiniProblems.ProblemPhyXMini0951.massReadout}
  \uses{def:physics:phyx-mini-0951:phyxminiproblems-problemphyxmini0951-massquantity}
  Read a physical mass in a selected unit
\end{definition}
\begin{proof}
  This is the declared model definition of mass Readout.
\end{proof}
\begin{definition}[length Readout]
  \label{def:physics:phyx-mini-0951:phyxminiproblems-problemphyxmini0951-lengthreadout}
  \lean{PhyXMiniProblems.ProblemPhyXMini0951.lengthReadout}
  \uses{def:physics:phyx-mini-0951:phyxminiproblems-problemphyxmini0951-lengthquantity}
  Read a physical length in a selected unit
\end{definition}
\begin{proof}
  This is the declared model definition of length Readout.
\end{proof}
\begin{definition}[acceleration In Meters Per Second Squared]
  \label{def:physics:phyx-mini-0951:phyxminiproblems-problemphyxmini0951-accelerationinmeterspersecondsquared}
  \lean{PhyXMiniProblems.ProblemPhyXMini0951.accelerationInMetersPerSecondSquared}
  \uses{def:physics:phyx-mini-0951:phyxminiproblems-problemphyxmini0951-accelerationmagnitude}
  Coherent-SI acceleration readout, in metres per second squared
\end{definition}
\begin{proof}
  This is the declared model definition of acceleration In Meters Per Second Squared.
\end{proof}
\begin{definition}[magnetic Flux Density In Teslas]
  \label{def:physics:phyx-mini-0951:phyxminiproblems-problemphyxmini0951-magneticfluxdensityinteslas}
  \lean{PhyXMiniProblems.ProblemPhyXMini0951.magneticFluxDensityInTeslas}
  \uses{def:physics:phyx-mini-0951:phyxminiproblems-problemphyxmini0951-magneticfluxdensitymagnitude}
  Coherent-SI magnetic-flux-density readout, in teslas
\end{definition}
\begin{proof}
  This is the declared model definition of magnetic Flux Density In Teslas.
\end{proof}
\begin{definition}[electric Current In Amperes]
  \label{def:physics:phyx-mini-0951:phyxminiproblems-problemphyxmini0951-electriccurrentinamperes}
  \lean{PhyXMiniProblems.ProblemPhyXMini0951.electricCurrentInAmperes}
  \uses{def:physics:phyx-mini-0951:phyxminiproblems-problemphyxmini0951-electriccurrentquantity}
  Coherent-SI signed-current readout, in amperes
\end{definition}
\begin{proof}
  This is the declared model definition of electric Current In Amperes.
\end{proof}
\begin{definition}[force In Newtons]
  \label{def:physics:phyx-mini-0951:phyxminiproblems-problemphyxmini0951-forceinnewtons}
  \lean{PhyXMiniProblems.ProblemPhyXMini0951.forceInNewtons}
  \uses{def:physics:phyx-mini-0951:phyxminiproblems-problemphyxmini0951-forcemagnitude}
  Coherent-SI force-magnitude readout, in newtons
\end{definition}
\begin{proof}
  This is the declared model definition of force In Newtons.
\end{proof}
\begin{definition}[torque In Newton Meters]
  \label{def:physics:phyx-mini-0951:phyxminiproblems-problemphyxmini0951-torqueinnewtonmeters}
  \lean{PhyXMiniProblems.ProblemPhyXMini0951.torqueInNewtonMeters}
  \uses{def:physics:phyx-mini-0951:phyxminiproblems-problemphyxmini0951-signedtorquequantity}
  Coherent-SI signed-torque readout, in newton metres
\end{definition}
\begin{proof}
  This is the declared model definition of torque In Newton Meters.
\end{proof}
\begin{definition}[degrees To Radians]
  \label{def:physics:phyx-mini-0951:phyxminiproblems-problemphyxmini0951-degreestoradians}
  \lean{PhyXMiniProblems.ProblemPhyXMini0951.degreesToRadians}
  Convert a dimensionless degree readout to radians
\end{definition}
\begin{proof}
  This is the declared model definition of degrees To Radians.
\end{proof}
\begin{definition}[Bar Point]
  \label{def:physics:phyx-mini-0951:phyxminiproblems-problemphyxmini0951-barpoint}
  \lean{PhyXMiniProblems.ProblemPhyXMini0951.BarPoint}
  The two labeled endpoints in the primary figure
\end{definition}
\begin{proof}
  This is the declared model definition of Bar Point.
\end{proof}
\begin{definition}[Cartesian Axis]
  \label{def:physics:phyx-mini-0951:phyxminiproblems-problemphyxmini0951-cartesianaxis}
  \lean{PhyXMiniProblems.ProblemPhyXMini0951.CartesianAxis}
  The two Cartesian axes drawn in the primary figure
\end{definition}
\begin{proof}
  This is the declared model definition of Cartesian Axis.
\end{proof}
\begin{definition}[In Plane Direction]
  \label{def:physics:phyx-mini-0951:phyxminiproblems-problemphyxmini0951-inplanedirection}
  \lean{PhyXMiniProblems.ProblemPhyXMini0951.InPlaneDirection}
  Directions in the plane of the supplied figure
\end{definition}
\begin{proof}
  This is the declared model definition of In Plane Direction.
\end{proof}
\begin{definition}[Page Normal Direction]
  \label{def:physics:phyx-mini-0951:phyxminiproblems-problemphyxmini0951-pagenormaldirection}
  \lean{PhyXMiniProblems.ProblemPhyXMini0951.PageNormalDirection}
  Directions perpendicular to the page
\end{definition}
\begin{proof}
  This is the declared model definition of Page Normal Direction.
\end{proof}
\begin{definition}[Bar Current Direction]
  \label{def:physics:phyx-mini-0951:phyxminiproblems-problemphyxmini0951-barcurrentdirection}
  \lean{PhyXMiniProblems.ProblemPhyXMini0951.BarCurrentDirection}
  Reference orientation for signed current along the bar
\end{definition}
\begin{proof}
  This is the declared model definition of Bar Current Direction.
\end{proof}
\begin{definition}[Torque Sense]
  \label{def:physics:phyx-mini-0951:phyxminiproblems-problemphyxmini0951-torquesense}
  \lean{PhyXMiniProblems.ProblemPhyXMini0951.TorqueSense}
  Sense of a torque in the page view
\end{definition}
\begin{proof}
  This is the declared model definition of Torque Sense.
\end{proof}
\begin{definition}[Pivoted Bar Figure]
  \label{def:physics:phyx-mini-0951:phyxminiproblems-problemphyxmini0951-pivotedbarfigure}
  \lean{PhyXMiniProblems.ProblemPhyXMini0951.PivotedBarFigure}
  \uses{def:physics:phyx-mini-0951:phyxminiproblems-problemphyxmini0951-barpoint, def:physics:phyx-mini-0951:phyxminiproblems-problemphyxmini0951-cartesianaxis}
  Literal qualitative content transcribed from image 951.png. The crosses encode a field into the page; the diagonal gold segment joins a to b; and the angle arc labeled theta is measured from the positive x axis at a
\end{definition}
\begin{proof}
  This is the declared model definition of Pivoted Bar Figure.
\end{proof}
\begin{definition}[Pivoted Conducting Bar Setup]
  \label{def:physics:phyx-mini-0951:phyxminiproblems-problemphyxmini0951-pivotedconductingbarsetup}
  \lean{PhyXMiniProblems.ProblemPhyXMini0951.PivotedConductingBarSetup}
  \uses{def:physics:phyx-mini-0951:phyxminiproblems-problemphyxmini0951-massquantity, def:physics:phyx-mini-0951:phyxminiproblems-problemphyxmini0951-lengthquantity, def:physics:phyx-mini-0951:phyxminiproblems-problemphyxmini0951-accelerationmagnitude, def:physics:phyx-mini-0951:phyxminiproblems-problemphyxmini0951-magneticfluxdensitymagnitude, def:physics:phyx-mini-0951:phyxminiproblems-problemphyxmini0951-electriccurrentquantity, def:physics:phyx-mini-0951:phyxminiproblems-problemphyxmini0951-forcemagnitude, def:physics:phyx-mini-0951:phyxminiproblems-problemphyxmini0951-signedtorquequantity, def:physics:phyx-mini-0951:phyxminiproblems-problemphyxmini0951-barpoint, def:physics:phyx-mini-0951:phyxminiproblems-problemphyxmini0951-inplanedirection, def:physics:phyx-mini-0951:phyxminiproblems-problemphyxmini0951-pagenormaldirection, def:physics:phyx-mini-0951:phyxminiproblems-problemphyxmini0951-barcurrentdirection, def:physics:phyx-mini-0951:phyxminiproblems-problemphyxmini0951-torquesense, def:physics:phyx-mini-0951:phyxminiproblems-problemphyxmini0951-pivotedbarfigure}
  Independent physical quantities describing the pivoted bar. The current and all three torques are observables, not definitions made from the requested answer
\end{definition}
\begin{proof}
  This is the declared model definition of Pivoted Conducting Bar Setup.
\end{proof}
\begin{definition}[Matches Problem And Figure Readouts]
  \label{def:physics:phyx-mini-0951:phyxminiproblems-problemphyxmini0951-matchesproblemandfigurereadouts}
  \lean{PhyXMiniProblems.ProblemPhyXMini0951.MatchesProblemAndFigureReadouts}
  \uses{def:physics:phyx-mini-0951:phyxminiproblems-problemphyxmini0951-massreadout, def:physics:phyx-mini-0951:phyxminiproblems-problemphyxmini0951-lengthreadout, def:physics:phyx-mini-0951:phyxminiproblems-problemphyxmini0951-accelerationinmeterspersecondsquared, def:physics:phyx-mini-0951:phyxminiproblems-problemphyxmini0951-magneticfluxdensityinteslas, def:physics:phyx-mini-0951:phyxminiproblems-problemphyxmini0951-degreestoradians, def:physics:phyx-mini-0951:phyxminiproblems-problemphyxmini0951-pivotedconductingbarsetup}
  Problem-statement and primary-raster readouts. The standard-gravity calibration makes explicit the conventional numerical constant needed to compare with the displayed answer choices. No current value appears here
\end{definition}
\begin{proof}
  This is the declared model definition of Matches Problem And Figure Readouts.
\end{proof}
\begin{definition}[Matches Pivoted Uniform Bar Scenario]
  \label{def:physics:phyx-mini-0951:phyxminiproblems-problemphyxmini0951-matchespivoteduniformbarscenario}
  \lean{PhyXMiniProblems.ProblemPhyXMini0951.MatchesPivotedUniformBarScenario}
  \uses{def:physics:phyx-mini-0951:phyxminiproblems-problemphyxmini0951-pivotedconductingbarsetup}
  Prose-level properties and signed-orientation conventions of the setup
\end{definition}
\begin{proof}
  This is the declared model definition of Matches Pivoted Uniform Bar Scenario.
\end{proof}
\begin{definition}[Has Physical Bar Parameters]
  \label{def:physics:phyx-mini-0951:phyxminiproblems-problemphyxmini0951-hasphysicalbarparameters}
  \lean{PhyXMiniProblems.ProblemPhyXMini0951.HasPhysicalBarParameters}
  \uses{def:physics:phyx-mini-0951:phyxminiproblems-problemphyxmini0951-massreadout, def:physics:phyx-mini-0951:phyxminiproblems-problemphyxmini0951-lengthreadout, def:physics:phyx-mini-0951:phyxminiproblems-problemphyxmini0951-accelerationinmeterspersecondsquared, def:physics:phyx-mini-0951:phyxminiproblems-problemphyxmini0951-magneticfluxdensityinteslas, def:physics:phyx-mini-0951:phyxminiproblems-problemphyxmini0951-pivotedconductingbarsetup}
  Positivity and non-vacuity conditions for the physical apparatus
\end{definition}
\begin{proof}
  This is the declared model definition of Has Physical Bar Parameters.
\end{proof}
\begin{definition}[Has Uniform Applied Magnetic Field]
  \label{def:physics:phyx-mini-0951:phyxminiproblems-problemphyxmini0951-hasuniformappliedmagneticfield}
  \lean{PhyXMiniProblems.ProblemPhyXMini0951.HasUniformAppliedMagneticField}
  \uses{def:physics:phyx-mini-0951:phyxminiproblems-problemphyxmini0951-magneticfluxdensityinteslas, def:physics:phyx-mini-0951:phyxminiproblems-problemphyxmini0951-pivotedconductingbarsetup}
  The scalar flux-density magnitude calibrates Physlib's spacetime-dependent magnetic vector field throughout the region marked by uniform crosses
\end{definition}
\begin{proof}
  This is the declared model definition of Has Uniform Applied Magnetic Field.
\end{proof}
\begin{definition}[Satisfies Uniform Bar Torque Laws]
  \label{def:physics:phyx-mini-0951:phyxminiproblems-problemphyxmini0951-satisfiesuniformbartorquelaws}
  \lean{PhyXMiniProblems.ProblemPhyXMini0951.SatisfiesUniformBarTorqueLaws}
  \uses{def:physics:phyx-mini-0951:phyxminiproblems-problemphyxmini0951-massreadout, def:physics:phyx-mini-0951:phyxminiproblems-problemphyxmini0951-lengthreadout, def:physics:phyx-mini-0951:phyxminiproblems-problemphyxmini0951-accelerationinmeterspersecondsquared, def:physics:phyx-mini-0951:phyxminiproblems-problemphyxmini0951-magneticfluxdensityinteslas, def:physics:phyx-mini-0951:phyxminiproblems-problemphyxmini0951-electriccurrentinamperes, def:physics:phyx-mini-0951:phyxminiproblems-problemphyxmini0951-forceinnewtons, def:physics:phyx-mini-0951:phyxminiproblems-problemphyxmini0951-torqueinnewtonmeters, def:physics:phyx-mini-0951:phyxminiproblems-problemphyxmini0951-pivotedconductingbarsetup}
  School-level statics laws for a uniform straight bar pivoted at one endpoint. The magnetic-torque equation is the integral of dτ = r × (I dℓ × B) along the bar. These are general relations among independent observables; they do not state the solved current
\end{definition}
\begin{proof}
  This is the declared model definition of Satisfies Uniform Bar Torque Laws.
\end{proof}
\begin{definition}[Is In Rotational Equilibrium]
  \label{def:physics:phyx-mini-0951:phyxminiproblems-problemphyxmini0951-isinrotationalequilibrium}
  \lean{PhyXMiniProblems.ProblemPhyXMini0951.IsInRotationalEquilibrium}
  \uses{def:physics:phyx-mini-0951:phyxminiproblems-problemphyxmini0951-torqueinnewtonmeters, def:physics:phyx-mini-0951:phyxminiproblems-problemphyxmini0951-pivotedconductingbarsetup}
  Rotational equilibrium means zero signed net torque about the pivot a
\end{definition}
\begin{proof}
  This is the declared model definition of Is In Rotational Equilibrium.
\end{proof}
\begin{definition}[Answer Choice]
  \label{def:physics:phyx-mini-0951:phyxminiproblems-problemphyxmini0951-answerchoice}
  \lean{PhyXMiniProblems.ProblemPhyXMini0951.AnswerChoice}
  Labels of the four current choices printed with the problem
\end{definition}
\begin{proof}
  This is the declared model definition of Answer Choice.
\end{proof}
\begin{definition}[answer Current In Amperes]
  \label{def:physics:phyx-mini-0951:phyxminiproblems-problemphyxmini0951-answercurrentinamperes}
  \lean{PhyXMiniProblems.ProblemPhyXMini0951.answerCurrentInAmperes}
  \uses{def:physics:phyx-mini-0951:phyxminiproblems-problemphyxmini0951-answerchoice}
  Current readout, in amperes, displayed beside each answer label
\end{definition}
\begin{proof}
  This is the declared model definition of answer Current In Amperes.
\end{proof}
\begin{definition}[recorded Dataset Answer]
  \label{def:physics:phyx-mini-0951:phyxminiproblems-problemphyxmini0951-recordeddatasetanswer}
  \lean{PhyXMiniProblems.ProblemPhyXMini0951.recordedDatasetAnswer}
  \uses{def:physics:phyx-mini-0951:phyxminiproblems-problemphyxmini0951-answerchoice}
  The answer label recorded by the source dataset
\end{definition}
\begin{proof}
  This is the declared model definition of recorded Dataset Answer.
\end{proof}
\begin{definition}[Answer Matches Current To Displayed Precision]
  \label{def:physics:phyx-mini-0951:phyxminiproblems-problemphyxmini0951-answermatchescurrenttodisplayedprecision}
  \lean{PhyXMiniProblems.ProblemPhyXMini0951.AnswerMatchesCurrentToDisplayedPrecision}
  \uses{def:physics:phyx-mini-0951:phyxminiproblems-problemphyxmini0951-electriccurrentinamperes, def:physics:phyx-mini-0951:phyxminiproblems-problemphyxmini0951-pivotedconductingbarsetup, def:physics:phyx-mini-0951:phyxminiproblems-problemphyxmini0951-answerchoice, def:physics:phyx-mini-0951:phyxminiproblems-problemphyxmini0951-answercurrentinamperes}
  Agreement to the displayed hundredth of an ampere. The half-unit in the last printed place is 0.005 A
\end{definition}
\begin{proof}
  This is the declared model definition of Answer Matches Current To Displayed Precision.
\end{proof}
\begin{definition}[Is Unique Closest Current Choice]
  \label{def:physics:phyx-mini-0951:phyxminiproblems-problemphyxmini0951-isuniqueclosestcurrentchoice}
  \lean{PhyXMiniProblems.ProblemPhyXMini0951.IsUniqueClosestCurrentChoice}
  \uses{def:physics:phyx-mini-0951:phyxminiproblems-problemphyxmini0951-electriccurrentinamperes, def:physics:phyx-mini-0951:phyxminiproblems-problemphyxmini0951-pivotedconductingbarsetup, def:physics:phyx-mini-0951:phyxminiproblems-problemphyxmini0951-answerchoice, def:physics:phyx-mini-0951:phyxminiproblems-problemphyxmini0951-answercurrentinamperes}
  A choice is strictly nearer to the required current than every rival
\end{definition}
\begin{proof}
  This is the declared model definition of Is Unique Closest Current Choice.
\end{proof}
\begin{lemma}[required Current formula]
  \label{lem:physics:phyx-mini-0951:phyxminiproblems-problemphyxmini0951-requiredcurrent-formula}
  \lean{PhyXMiniProblems.ProblemPhyXMini0951.requiredCurrent_formula}
  \uses{def:physics:phyx-mini-0951:phyxminiproblems-problemphyxmini0951-massreadout, def:physics:phyx-mini-0951:phyxminiproblems-problemphyxmini0951-lengthreadout, def:physics:phyx-mini-0951:phyxminiproblems-problemphyxmini0951-accelerationinmeterspersecondsquared, def:physics:phyx-mini-0951:phyxminiproblems-problemphyxmini0951-magneticfluxdensityinteslas, def:physics:phyx-mini-0951:phyxminiproblems-problemphyxmini0951-electriccurrentinamperes, def:physics:phyx-mini-0951:phyxminiproblems-problemphyxmini0951-pivotedconductingbarsetup, def:physics:phyx-mini-0951:phyxminiproblems-problemphyxmini0951-matchespivoteduniformbarscenario, def:physics:phyx-mini-0951:phyxminiproblems-problemphyxmini0951-hasphysicalbarparameters, def:physics:phyx-mini-0951:phyxminiproblems-problemphyxmini0951-satisfiesuniformbartorquelaws, def:physics:phyx-mini-0951:phyxminiproblems-problemphyxmini0951-isinrotationalequilibrium}
  Torque balance gives the general required-current relation I = m g cos(theta) / (B L). This is a conclusion, not a premise field
\end{lemma}
\begin{proof}
  The conclusion follows from the governing relations and figure readouts named in the statement of required Current formula.
\end{proof}
% --- Archon declaration topology end ---
% --- Archon physics formalization source end ---
